
les $\xi_\ell$ \ill{} dépendant du choix d'un pt géom.\ de $X$, mais à
isom.\ unique près (dépendance du $\pi_1(X)$) [$\pi_1(X)$ opère sur
$\xi_\ell$] \struck{sont} indépendants de ce choix.
\note{« dépendance du $\pi_1(X)$ » est souligné deux fois}

Les $\xi_p^{\alpha}$ dépendent du choix d'un \struck{corps} pt géom.\ [de
$X$] à valeurs dans un corps parfait sans plus, si on veut les
$k(x)^{\mathrm{pf}}$ (clôture parfaite de $k(x)$), et pour $k$ donné, il ne
semble pas qu'il y ait de choix des isom.\ entre les $\xi_{p,a}^{(\alpha)}$
correspondant aux divers pts à valeurs dans $k$ (c'est-à-dire « points
\ill{} » du \ill{} géom.\ qui \ill{} des \uncertain{maximaux}), s'ils sont
localisés en le même $x \in X$, \struck{\ill{}} \struck{\ill{}} si $k$ est
alg.\ clos, ils sont conjugués par autom.\ de $k$ [(éléments de
$W_p(k)$)] \ill{} alg., mais sont \ill{} conjugués \ill{} autom.\ de $k$
… il est à \ill{} général qu'il semble \ill{} l'un \ill{} de
l'autre… Le gpe $\pi_1(X)$ n'opère plus sur $\xi_p^{\alpha}$,
\marginal{En revanche, on a comme structure supplémentaire les opérations
$F$, $V$ avec $FV = p$}
par contre le groupe des autom.\ de $k$ sur $k(x)$
\struck{\ill{}} opère ½ linéairement sur $\xi_p^{(\alpha)}$ ---
\ill{} fibre $\ell$-adique en \ill{}
\note{« ½ linéairement » est ajouté dans l'interligne et souligné ; une
insertion biffée au-dessus de la ligne se termine par un signe de renvoi qui
mène à la note marginale ci-dessous}

\marginal{\ill{} élégante. D'autre part, un autom.\ $\sigma$ de $k$ qui
induit sur $k(x)$ l'identité transforme en général $\xi_{x,p}$ en un
$\xi_{x,p}^{\sigma}$ qui n'est plus isom.\ à $\xi_{x,p}$ ---}
\note{les deux notes marginales sont écrites verticalement dans la marge
gauche ; dans la seconde, les indices et exposants des deux premiers $\xi$
sont douteux}

Notations,
\begin{enumerate}
  \item[(i)] $\xi_{\bar{x},\ell}$ \ill{} ; $\xi_{x,\ell}$ \ill{} pour
  $x \in X$ ou un pt de $x$ à val.\ dans un corps quelc.\ ($\xi_{x,\ell}$
  déf.\ à isom.\ non unique près)
  \item[(ii)] $\xi_{x,p}$, $\xi_{x,p}^{0}$ : $x$ pt de $X$ à val.\ dans un
  corps \uncertain{parfait}.
\end{enumerate}

\marginal{\uncertain{Gros} point \uncertain{subtil} \ill{} \ill{}}
\note{note de trois lignes soulignées, en bas de la marge gauche}

